% This file archives material that is not included in the paper.

\iffalse
\subsubsection{An exact configuration LP for the rounded problem}

For $S\in\widehat\cC$, let $
    q_S=\frac{1}{2}c\br{\delta_E(S)}$ be the cost of cutting the edges in $\delta_E(S)$.
Note, that the factor $1/2$ compensates for counting each supply edge between two parts in both boundaries. Introduce one variable $\lambda_S$ for each rounded configuration:
\begingroup
\renewcommand{\theHequation}{tree-config-primal}
\begin{equation}\label{lp:tree-config-primal}
\tag{$P_{\widehat d,L}$}
\begin{aligned}
    \min\quad
        & \sum_{S\in\widehat\cC}q_S\cdot \lambda_S \\
    \text{subject to}\quad
        & \sum_{S\in\widehat\cC\colon v\in S}\lambda_S=1
            && \forall v\in V,\\
        & \lambda_S\geq0
            && \forall S\in\widehat\cC.
\end{aligned}
\end{equation}
\endgroup
Its dual has one unrestricted variable $\alpha_v$ for each vertex:
\begingroup
\renewcommand{\theHequation}{tree-config-dual}
\begin{equation}\label{lp:tree-config-dual}
\tag{$D_{\widehat d,L}$}
\begin{aligned}
    \max\quad
        & \sum_{v\in V}\alpha_v \\
    \text{subject to}\quad
        & \sum_{v\in S}\alpha_v\leq q_S
            && \forall S\in\widehat\cC,\\
        & \alpha_v\in\mathbb{R}
            && \forall v\in V.
\end{aligned}
\end{equation}
\endgroup

For a dual vector $\alpha$, let $
    \alpha(S)=\sum_{v\in S}\alpha_v.$
We define the \textsc{Pricing Problem} as
\begin{equation}\label{eq:tree-pricing}
    \max_{S\in\widehat\cC}\brc{\alpha(S)-q_S}
\end{equation}
In order to solve the dual, using the ellipsoid method, we need to provide a separation oracle for the constraints of type~\eqref{lp:tree-config-dual}. It is easy to see, that this is equivalent to solving the pricing problem~\eqref{eq:tree-pricing}. It what follows we show an exact dynamic program for the pricing problem with runtime polynomial in $L$ and $n$.

For each vertex, let $d_c(v)=\sum_{e\ni v}c_e$ and $\beta_v=\alpha_v-\frac{1}{2}d_c(v).$
For every connected $S\subseteq V$, we can now rewrite
\begin{equation}\label{eq:tree-pricing-profit}
    \alpha(S)-q_S
    =\sum_{v\in S}\beta_v+c\br{E(S)}.
\end{equation}

Root $T$ arbitrarily. For a vertex $v$, let $T_v$ be its descendant subtree, and let $P_v(\ell)$ be the maximum value of the right-hand side of~\eqref{eq:tree-pricing-profit} over connected sets $S\subseteq V(T_v)$ containing $v$ and satisfying $\sum_{u\in S}\widehat d(u)=\ell$. An impossible state has value $-\infty$.

If the children of $v$ are $u_1,\ldots,u_k$, denote by
$P_v^{(j)}(\ell)$ the maximum value of the right-hand side of~\eqref{eq:tree-pricing-profit} over connected sets $S\subseteq V(T_v)$ containing $v$ and the first $j$ children and satisfying $\sum_{u\in S}\widehat d(u)=\ell$.
Initialize
\[
    P_v^{(0)}(\ell)
    =\begin{cases}
        \beta_v,&\ell=\widehat d(v),\\
        -\infty,&\text{otherwise}.
    \end{cases}
\]
After processing the first $j-1$ children, merge the table of $u_j$ by
\begin{align*}
    P_v^{(j)}(\ell)
    =\max\brc{
        P_v^{(j-1)}(\ell),
        \max_{0\leq a,b\leq L\colon a+b=\ell}
        \brc{P_v^{(j-1)}(a)+P_{u_j}(b)+c(vu_j)}
    }.
\end{align*}
Here, the first term corresponds to not including the $j$-th child in $S$, and the second term corresponds to including it.

After all children are processed, set $P_v=P_v^{(k)}$. Every nonempty connected set has a unique vertex closest to the root, so the optimum in~\eqref{eq:tree-pricing} equals $
    \max_{v\in V, 0\leq\ell\leq L}\brc{P_v(\ell)}.$ 
Since the time required to calculate one value of $P_v^{(j)}(\ell)$ is polynomial $O\br{L}$ and there are $O(nL)$ such values for all possible $v$ and $j$, the overall time complexity is $O\br{nL^2}$ as required.
By Lemma~\ref{lem:rounded-degree-family}, this is polynomial in the input size and in $1/(2\cdot\eta-1-\rho)$.

For a vector $\lambda$ indexed by $\widehat\cC$, define
\[
    \supp(\lambda)=\brc{S\in\widehat\cC:\lambda_S>0}.
\]
to be the support of the vector $\lambda$.

\begin{lemma}\label{lem:tree-config-primal-recovery}
    In polynomial time one can compute an optimal solution $\lambda^*$ of~\eqref{lp:tree-config-primal} in polynomial time with $\supp(\lambda^*)$ of polynomial size.
\end{lemma}
\begin{proof}
    Run the rational ellipsoid optimization-separation algorithm for~\eqref{lp:tree-config-dual} in its primal-recovery form. Whenever the dynamic program finds a violated constraint, store its configuration. The ellipsoid algorithm makes polynomially many calls and stores a polynomial-size family $\cR\subseteq\widehat\cC$. Primal recovery states that the LP restricted to $\cR$ has the same optimum as the full primal. Solving that restricted LP and implcitly extending its solution by zero outside $\cR$ proves the claim.
\end{proof}

\begin{theorem}\label{thm:tree-config-lossless-rounding}
    Let $\lambda$ be a feasible solution of~\eqref{lp:tree-config-primal} with value $
        z=\sum_{S\in\widehat\cC}q_S\cdot\lambda_S.
    $
    In time polynomial in the size of $\supp(\lambda)$, one can find a partition $\cP\subseteq\widehat\cC$ of $V$ such that
    \[
        c\br{\delta_E(\cP)}
        =\sum_{S\in\cP}q_S
        \leq z.
    \]
\end{theorem}
\begin{proof}
    Let $\cF=\supp(\lambda)$. Construct the intersection graph $J$ whose vertices are the configurations in $\cF$, with two configurations adjacent if and only if they intersect.

    Since the members of $\cF$ are subtrees of $T$, the graph $J$ is chordal. Moreover, subtrees of a tree have the Helly property: every pairwise-intersecting family of subtrees has a common vertex. Therefore, every clique $\cK$ of $J$ has a vertex $v\in V$ common to all its configurations, and hence
    \[
        \sum_{S\in\cK}\lambda_S
        \leq
        \sum_{\substack{S\in\cF\\v\in S}}\lambda_S
        =1.
    \]
    Recall that a graph is perfect if, in every induced subgraph, the chromatic number equals the size of a largest clique. Chordal graphs are perfect. For a graph $J$, its independent-set polytope is the convex hull of the incidence vectors of all independent sets of $J$. A standard polyhedral characterization of perfect graphs states that this polytope is exactly
    \[
        \left\{
            x\in\mathbb{R}_{\geq0}^{\cF}:
            \sum_{S\in\cK}x_S\leq1
            \text{ for every clique $\cK$ of $J$}
        \right\}.
    \]
    The displayed inequalities are called the clique inequalities: an independent set contains at most one vertex from each clique. The vector $\lambda$, restricted to $\cF$, is non-negative and satisfies every clique inequality by the argument above. It therefore belongs to the independent-set polytope of $J$. Equivalently, it is a convex combination of incidence vectors of independent sets of $J$.

    We now give an explicit rounding procedure. Let $
        Q=\sum_{S\in\cF}q_S$ and $M=Q+1$. Now, give each vertex $S$ of $J$ the weight $
        p_S=M\spr{S}-q_S.$
    Compute a maximum-weight independent set $\cI$ of the chordal graph $J$. This can be done in polynomial time. Since $\lambda$ lies in the independent-set polytope of $J$,
    \begin{align*}
        \sum_{S\in\cI}p_S
        &\geq\sum_{S\in\cF}p_S\lambda_S\\
        &=M\sum_{S\in\cF}\spr{S}\lambda_S-z\\
        &=M\sum_{v\in V}\sum_{\substack{S\in\cF\\v\in S}}\lambda_S-z\\
        &=M\spr{V}-z.
    \end{align*}
    For every $S\in\cF$, feasibility of $\lambda$ implies $\lambda_S\leq1$. Consequently, $
        z\leq Q<M.$ and therefore it follows that $
        M\spr{V}-z>M\br{\spr{V}-1}.$

    The configurations in an independent set are pairwise disjoint. Thus, if $\cI$ failed to cover some vertex of $V$, then
    \[
        \sum_{S\in\cI}p_S
        =M\spr{\bigcup_{S\in\cI}S}-\sum_{S\in\cI}q_S
        \leq M\br{\spr{V}-1},
    \]
    which is a contradiction. Hence, $\cI$ is a partition of $V$. Finally,
    \[
        M\spr{V}-\sum_{S\in\cI}q_S
        =\sum_{S\in\cI}p_S
        \geq M\spr{V}-z,
    \]
    and therefore $ \sum_{S\in\cI}q_S\leq z.$ Since every edge of $\delta_E(\cI)$ belongs to the boundaries of exactly two parts, the definition of $q_S$ gives
    \[
        c\br{\delta_E(\cI)}
        =\sum_{S\in\cI}q_S
        \leq z.
    \]
    Taking $\cP=\cI$ proves the theorem.
\end{proof}

\begin{theorem}\label{thm:tree-degree-config-direct}
    Let $0<\rho<1$ and $\eta>(1+\rho)/2$ be rational. For arbitrary demand-edge weights on a supply tree, there is an algorithm, polynomial in the input size and in $
        \frac{1}{2\cdot\eta-1-\rho}$,
    that returns an $\eta$-feasible solution $F$ with
    \[
        c\br{F}\leq\OPTdem{T}{\rho}.
    \]
\end{theorem}
\begin{proof}
    Round the demand degrees as in~\eqref{eq:rounded-degree-definition} before constructing any LP. Solve the configuration LP over the resulting family $\widehat\cC$ using the exact separation oracle above, recover a polynomial-support optimum, and apply Theorem~\ref{thm:tree-config-lossless-rounding}. By Lemma~\ref{lem:rounded-degree-family}, every part returned by the rounding is $\eta$-feasible. The same lemma shows that every part of an optimal $\rho$-feasible partition belongs to $\widehat\cC$. Therefore, the rounded configuration-LP optimum, and hence the cost returned by lossless rounding, is at most $\OPTdem{T}{\rho}$.
\end{proof}
\fi
